\documentclass{article}
\usepackage{graphicx}
\usepackage[margin=1.5cm]{geometry}
\usepackage{amsmath}

\begin{document}
\twocolumn

\title{Wednesday Warm Up: Unit 0, Review of 150}
\author{Prof. Jordan C. Hanson}

\maketitle

\section{Review of 150}

\begin{enumerate}
\item Suppose a physical therapy patient is asked to shove a medicine ball forward off the edge of a table to help rebuild the strength of their shoulders.  The medicine ball weighs 7 kg.  (a) If the patient is able to give the ball a speed of 1 m s$^{-1}$, what is the momentum of the ball? (b) Suppose the medicine ball has a diameter of 10 cm.  The moment of inertia for a solid sphere is $I = \frac{2}{5} m r^2$.  What is the angular momentum of the ball if it is spun at 1 rotation per second? (c) What is the rotational kinetic energy? \\ \vspace{5cm}
\end{enumerate}
\section{Mathematics Review}
\begin{enumerate}
\item Compute the dot-product, $\vec{a} \cdot \vec{b}$.
\begin{itemize}
\item $\vec{a} = -2\hat{i} + 2\hat{j}$, $\vec{b} = 4\hat{i} + 4\hat{j}$.
\item $\vec{a} = 2\hat{i} - 1\hat{j}$, $\vec{b} = -2\hat{i} - 3\hat{j}$.
\end{itemize}
\item Compute the cross-product, $\vec{a} \times \vec{b}$.
\begin{itemize}
\item $\vec{a} = -2\hat{i} + 2\hat{j}$, $\vec{b} = 4\hat{i} + 4\hat{j}$.
\item $\vec{a} = 2\hat{i} - 1\hat{j}$, $\vec{b} = -2\hat{i} - 3\hat{j}$.
\end{itemize}
\end{enumerate} \vspace{4cm}
\section{Unit 0: Electrostatics I}
\begin{enumerate}
\item The SI unit of \textbf{charge} is called a \textit{Coulomb} (1 C).  A Coulomb is the amount of charge that flows through a wire in one second, given that the \textit{current} is 1 \textit{Amp}. Define current as $I = \Delta Q/\Delta t$, where $Q$ is the charge as a function of time, and $\Delta Q$ is the change in the charge.  (a)  Suppose we are charging a car battery with an average current of $I = \Delta Q/\Delta t = 2$ Amps, for 5 hours.  How many Coulombs of charge did we put in the battery? (b) If a battery has 7 Amp hr of charge, for how long can it provide 0.5 Amps of current? \\ \vspace{4cm}
\item If $1.80 \times 10^{20}$ electrons move through a pocket calculator during a full day's operation, how many coulombs of charge moved through it?\footnote{One electron has $1.6 \times 10^{-19}$ C of charge.}
\end{enumerate}

\end{document}
